\begin{figure}[htbp]
\centering
\begin{tikzpicture}[tablero A]
 \begin{scope}
  \cuadriculaA{4}\trazarRejillaA{4}
  \ataqueA{3}{2}{2}{1}\aspaA{2}{1}
  \aroA{0}{0}\reinaA[colorAro]{0}{0}\reinaA{1}{3}\reinaA{3}{2}
  \panelA{2,4.38}{(a) Remove one queen}
 \end{scope}
 \draw[-{Stealth[length=2mm]},tintaReina,line width=.9pt] (4.9,2)--(6.65,2);
 \begin{scope}[shift={(8,0)}]
  \cuadriculaA{4}\trazarRejillaA{4}
  \ataqueA{3}{2}{2}{1}\aspaA{2}{1}
  \reinaA{1}{3}\reinaA{3}{2}
  \aroA{0}{1}\reinaA[colorAro]{0}{1}
  \aroA{2}{0}\reinaA[colorAro]{2}{0}
  \panelA{2,4.38}{(b) Place two queens}
 \end{scope}
\end{tikzpicture}
\caption{A legal repair takes a placement from $\Omega_1(4)$ to $\Omega_0(4)$.
The rings mark the pieces involved. The cross marks the same square $(2,1)$ in both panels.
Reading from right to left recovers the reverse operation with the ordered pair $(2,0),(0,1)$.}
\label{fig:reparacion-tablero}
\end{figure}
